\par\bigskip
\begin{center}
\pdfbookmark[2]{Gauss an Olbers, 30. November 1810}{pallas-brief-18101130}
\textsc{Gauss} an \textsc{Olbers}, 30. November 1810.
\end{center}
\noindent\hspace*{\parindent}Ich kann nicht umhin, Ihnen wenigstens mit ein paar Worten zu schreiben, dass ich jetzt beinahe fertig bin mit der Arbeit, über welche ich Ihnen neulich schrieb, und dass sie, wenn mich nicht alles trügt, über alle meine Erwartung befriedigend ausfallen wird. Gestern berechnete ich zuerst aus den 4 ersten Oppositionen die Normalelemente, die mit den vorher integrirten Störungen verbunden jene genau darstellten, und diese wichen von der 5. Opposition \(4'\), von der 6. \(20'\) zu meinem grossen Schrecken ab. Ich entdeckte hiedurch, dass in meine ganze Rechnung ein Fehler eingeschlichen war, indem der Haupttheil der Störung der mittlern Bewegung mit falschem Zeichen genommen war. Heute habe ich nun alles redressirt, und nach den vorläufig leni calamo nur mit den kleinen Tafeln berechneten Normalelementen werden diese
\begin{center}\begingroup\small\setlength{\tabcolsep}{4pt}
\begin{tabular}{lrrrrrr}
bei den 6 Längen & \(8''\) & \(1''\) & \(8''\) & \(0''\) & \(6''\) & \(5''\)\\
bei den 6 Breiten & \(2''\) & \(5''\) & \(6''\) & \(5''\) & \(17''\) & \(3''\)\\
\end{tabular}\endgroup\end{center}
abweichen. Die feinere Politur muss erst noch hinzukommen (die Breite der 5. Opposition war bekanntlich schlecht beobachtet, und ich traute ihr einen Fehler von \(30''\) zu). Die rein elliptischen Elemente, die am \textls[180]{besten} übereinstimmten, differirten
\begin{center}\begingroup\small\setlength{\tabcolsep}{4pt}
\begin{tabular}{lrrrrrr}
in der Länge & \(-111''\) & \(+59''\) & \(+20''\) & \(+86''\) & \(+137''\) & \(-217''\)\\
in der Breite & \(-8''\) & \(-37''\) & \(0''\) & \(+25''\) & \(+29''\) & \(+83''.\)\\
\end{tabular}\endgroup\end{center}
\par
\noindent\hspace*{\parindent}Um Ihnen einige Ideen von den Störungen der Elemente zu geben, setze ich ein paar her:
\begin{center}\begingroup\scriptsize\setlength{\tabcolsep}{2pt}\renewcommand{\arraystretch}{1.2}
\begin{tabular}{lccccccc}
{\AstroSymbols ☍}\ .\,.\,.\,.\,.\,.\,.\,.\,.\,.\,.\,. & 1803 & 1804 & 1805 & 1807 & 1808 & 1809 & 1811\\
\(\left.\begin{gathered}\text{Motus diurnus}\\\text{medius sider.}\end{gathered}\right\}\) & \(770''{,}721\) & \(770''{,}389\) & \(770''{,}133\) & \(770''{,}497\) & \(770''{,}154\) & \(768''{,}293\) & \(768''{,}860\)\\
\(\varphi\)\ .\,.\,.\,.\,.\,.\,.\,.\,.\,.\,.\,. & \(14^{0}12'47''\) & \(11'33''\) & \(10'59''\) & \(10'10''\) & \(8'10''\) & \(1'35''\) & \(13^{0}58'57''.\)\\
\end{tabular}\endgroup\end{center}
\par
